% Source: https://artofproblemsolving.com/community/c3h2666421_question_about_square_roots
% Stable identity: aops-topic-2666421-q01
Suppose $\sqrt{23+\sqrt5}=a+\sqrt b$ with $a,b$ rational. Squaring gives $a^2+b+2a\sqrt b=23+\sqrt5$, so
\[2a\sqrt b=\sqrt5+r,\qquad\text{where }r=23-a^2-b\text{ is rational}.\]
Squaring again, $4a^2b=5+r^2+2r\sqrt5$. The left side and $5+r^2$ are rational while $\sqrt5$ is irrational, so $r=0$. Hence
\[a^2+b=23\qquad\text{and}\qquad4a^2b=5.\]
Put $t=a^2$; then $b=23-t$ and $4t(23-t)=5$, i.e. $4t^2-92t+5=0$. Its roots are
\[t=\dfrac{92\pm\sqrt{92^2-80}}{8}=\dfrac{92\pm\sqrt{8384}}{8}=\dfrac{23\pm2\sqrt{131}}{2},\]
using $8384=64\cdot131$. Since $131$ is not a perfect square, both roots are irrational, but $t=a^2$ is rational. This contradiction shows that no such $a,b$ exist. $\blacksquare$

(Remark: in general $\sqrt{A+\sqrt B}$ with rational $A,B$ denests as $\sqrt x+\sqrt y$ with rational $x,y$ exactly when $A^2-B$ is the square of a rational number; here $23^2-5=524$ is not.)
